% ============================================================
% SF-3: The Quark Sector from 600-Cell Geometry
%       Masses, Strong Coupling, Koide Phase, and Generation Count
%       from a Single Calibration (m_e)
% Conscious Point Physics Flagship Paper Series (SF-line)
% Version history: flagship_papers/quarks/documentation_suite/changelog-sf-3.md
% ============================================================

\documentclass[11pt,letterpaper]{article}
\usepackage[utf8]{inputenc}
\usepackage[margin=1in]{geometry}
\usepackage[T1]{fontenc}
\usepackage{amsmath,amssymb,amsfonts,amsthm}
\usepackage{graphicx}
\usepackage{xcolor}
\usepackage{hyperref}
\usepackage{booktabs}
\usepackage{array}
\usepackage{enumitem}
\usepackage{float}
\usepackage[font=small, labelfont=bf]{caption}
\usepackage{parskip}
\usepackage{mdframed}

\hypersetup{
    colorlinks=true,
    linkcolor=blue,
    citecolor=blue,
    urlcolor=blue
}

\newtheorem{theorem}{Theorem}[section]
\newtheorem{proposition}[theorem]{Proposition}
\newtheorem{lemma}[theorem]{Lemma}
\newtheorem{corollary}[theorem]{Corollary}
\newtheorem{conjecture}[theorem]{Conjecture}
\newtheorem{definition}[theorem]{Definition}
\newtheorem{remark}[theorem]{Remark}
\newtheorem{openproblem}{Open Problem}

% --- Standard CPP macros ---
\newcommand{\phig}{\varphi}            % golden ratio
\newcommand{\Tr}{\operatorname{Tr}}
\newcommand{\SSV}{\mathrm{SSV}}
\newcommand{\DP}{\mathrm{DP}}
\newcommand{\CP}{\mathrm{CP}}

% --- Paper-specific macros ---
\newcommand{\Mzero}{M_0}
\newcommand{\sinsqthetaW}{\sin^2\theta_W}
\newcommand{\alphas}{\alpha_s}
\newcommand{\CF}{C_F}
\newcommand{\Kthree}{\mathrm{K}_3}
\newcommand{\CKM}{\mathrm{CKM}}
\newcommand{\PMNS}{\mathrm{PMNS}}
\newcommand{\mc}{m_c}
\newcommand{\me}{m_e}
\newcommand{\dgr}{^{\circ}}

\title{\textbf{SF-3: The Quark Sector from 600-Cell Geometry}\\[6pt]
{\large Masses, Strong Coupling, Koide Phase, and Generation Count from a Single Calibration}\\[4pt]
{\large Conscious Point Physics Flagship Paper Series}\\[4pt]
{\Large \textbf{Version 1.2 SHIPPED --- 14 June 2026}}}

\author{Thomas Lee Abshier, ND \and Claude Opus (Anthropic)}

\date{\normalsize Hyperphysics Institute, Kalispell, Montana\\
\url{https://hyperphysics.com}\\
\href{mailto:drthomas007@protonmail.com}{drthomas007@protonmail.com}}

\begin{document}
\maketitle

% ============================================================
% ABSTRACT
% ============================================================
\begin{abstract}
\noindent\textbf{Paper type:} Flagship synthesis paper. Establishes that the heavy-quark mass spectrum (strange, charm, bottom, top), the strong coupling, the quark Koide phase, and the three-generation count all follow from a \emph{single} calibration (the electron mass $\me$) plus 600-cell geometry and $SU(3)$ colour, with zero shape parameters fitted.

The Standard Model treats the six quark masses, the strong coupling $\alphas$, the number of generations, and the Cabibbo--Kobayashi--Maskawa ($\CKM$) mixing matrix as independent measured inputs. SF-3 assembles results from the shipped Conscious Point Physics (CPP) Standard-Model series (SM-7, SM-8, SM-9, SM-10) and Strong-Sector series (SS-1, SS-2) into a unified account in which (i) the heavy-quark masses $m_s, m_c, m_b, m_t$ derive from the zero-parameter cage formula $M_q = \me\,(z/\phig)\,V^{7/3}$ (with a top-relay multiplier $z\CF$), reproducing all four masses across four orders of magnitude at RMS $2.1\%$ from $\me$, the 600-cell coordination $z=12$, the golden ratio $\phig$, and the colour Casimir $\CF=4/3$ alone; (ii) the strong coupling $\alphas = 5/(8\phig) \approx 0.386$ emerges as the \emph{face-mode} fraction of the same 600-cell spectral trace that yields the Weinberg angle as its \emph{edge-mode} fraction $\sinsqthetaW = 3/(8\phig)$, with the exact complementarity $\sinsqthetaW + \alphas = 1/\phig$; (iii) the quark Koide phase $\theta_{\mathrm{quark}} = 124.04\dgr$ (vs PDG $124.09\dgr$, $0.05\%$) derives from the net isotropic self-energy shift built entirely from $\alphas$, $\sinsqthetaW$, and $z$; and (iv) the count of exactly three generations, with no fourth quark, is selected within the SM-8 antipodal-identification model.

\textbf{Calibration resolution.} Because the SM-8 cage formula \emph{predicts} the charm mass ($m_c \approx 1249$~MeV, $-1.6\%$) rather than taking it as input, the second calibration ($m_c$) that appeared in the earlier SM-7 quark-Koide route is \emph{redundant}: $m_c$ is demoted from a calibration to a derived quantity. This restores a uniform single-$\me$ calibration headline across charged leptons (SF-1), neutrinos (SF-4), and quarks (SF-3) --- the spine of the SF-line's ``hierarchy without hierarchy'' argument. A worked separation shows that the quark Koide phase depends on $\alphas$, $\sinsqthetaW$, and $z$ but \emph{not} on the mass amplitude or on $m_c$; no re-grounding of the phase machinery on a derived $m_c$ is required.

\textbf{The honestly-open piece.} SF-3 does \emph{not} derive the $\CKM$ mixing matrix or the quark $CP$ phase. This is structurally parallel to SF-4's open neutrino $\delta_{CP}$ in the limited sense that both flagships derive masses while deferring a mixing-sector $CP$ observable (a parallel of posture, not an equivalence of difficulty or status): SF-3 ships the quark masses and generation count at zero parameters with the mixing sector deferred (registered as OPEN-FP-3-CKM). The uniform ``masses derived, mixing-sector open'' posture across the fermion flagships is stated explicitly. The paper closes with four falsifiers, including a fourth-generation quark and a deviation of the measured quark Koide phase from $124.04\dgr$.
\end{abstract}

\medskip\noindent
\textbf{Keywords:} quark masses, strong coupling, $\alpha_s$, Koide formula, quark Koide phase, 600-cell lattice, golden ratio, electroweak--strong mode complementarity, three generations, fourth-generation falsifier, CKM matrix, Conscious Point Physics, single-$m_e$ calibration, zero-parameter prediction, Standard Model emergence, flagship paper.

\bigskip\noindent
\textbf{Plain Language Summary:}
The six quarks span five orders of magnitude in mass --- from the few-MeV up and down quarks to the 173-GeV top --- and the Standard Model simply measures these values without explaining them. This paper proposes that the four heavy quark masses (strange, charm, bottom, top) are fixed by the geometry of a four-dimensional lattice called the 600-cell, using only one calibration: the electron's mass. The same geometry that sets the electroweak ``Weinberg angle'' also sets the strength of the strong nuclear force, and the two add up to a single golden-ratio number exactly. A subtle angle that organizes the quark masses (the ``Koide phase'') comes out to within $0.05\%$ of its measured value. The lattice also forces there to be exactly three families of quarks and no fourth. The one thing the paper does \emph{not} explain is how quarks of different families mix --- the ``CKM'' matrix --- which it flags as honestly unfinished, exactly as the companion neutrino paper flags its own mixing parameter as unfinished.

\bigskip
\raggedright
\tableofcontents
\newpage

% ============================================================
% SECTION 1: Introduction
% ============================================================
\section{Introduction}

The quark sector poses four quantitative questions that the Standard Model (SM) answers only by measurement:

\begin{enumerate}[leftmargin=2em]
\item \textbf{The mass spectrum.} Why do the six quark masses take the values they do, spanning roughly five orders of magnitude from the light up/down quarks to the top?
\item \textbf{The strong coupling.} Why is $\alphas$ what it is, and is its value related to the electroweak couplings?
\item \textbf{The generation count.} Why exactly three generations, and is a fourth excluded?
\item \textbf{Mixing.} What sets the $\CKM$ matrix elements and the quark $CP$-violating phase?
\end{enumerate}

This paper assembles a Conscious Point Physics (CPP) account that answers the first three from a single calibration and the 600-cell substrate, and that locates the fourth as an explicitly open problem. None of the physics here is new derivation: every load-bearing result is inherited from shipped CPP papers and reframed into a single coherent quark-sector narrative. The contribution of SF-3 is the \emph{synthesis} --- including the adjudication of a calibration question that had been left open across the source corpus --- and the placement of the quark sector within the SF-line grand-unification program (SF-7).

\begin{mdframed}
\textbf{What SF-3 does \emph{not} claim.} SF-3 introduces no new derivation; every quantitative result is inherited from SM-7/8/9/10 and SS-1/2. In particular: (i) $\alphas$ and $\sinsqthetaW$ are \emph{structural correspondences} --- mode fractions of the 600-cell adjacency spectrum that match the measured couplings --- not gauge couplings derived from renormalization-group running; (ii) the three-generation count is \emph{selected within} the SM-8 antipodal-identification model, an axiom-conditional geometric exclusion, not a dynamical no-go theorem; (iii) the $\CKM$ mixing matrix and the quark $CP$-violating phase are \emph{not} derived (registered as OPEN-FP-3-CKM); and (iv) SM-10's first-principles cascade remains a calibrated model --- the mass \emph{values} do not depend on it. Proposition~\ref{prop:phaseindep} is a bookkeeping separation, not a theorem.
\end{mdframed}

CPP derives Standard-Model structure from the geometry of the 600-cell polytope (vertices $V=120$, edges $E=720$, faces $F=1200$, cells $C=600$, vertex coordination $z=12$). Conscious Points (CPs) occupy lattice sites and exchange displacement-increment information along Space-Stress-Vector ($\SSV$) gradients; particles are geometric cage structures, and the golden ratio $\phig=(1+\sqrt5)/2$ recurs because it is built into the 600-cell metric. The quark sector draws on two shipped derivation chains: the Standard-Model emergence series (SM-7 through SM-10), which supplies the mass formula and the strong coupling, and the Strong-Sector series (SS-1, SS-2), which supplies $SU(3)$ colour and the colour Casimir $\CF=4/3$.

\subsection{Open Problems Addressed}
This paper consolidates and reframes results bearing on the quark mass and coupling sector and registers one new open problem:
\begin{itemize}[leftmargin=2em]
\item It \emph{resolves} the single-$\me$ calibration question for the quark sector (whether the heavy-quark spectrum requires a second calibration $m_c$) in favour of a single calibration, by adopting the SM-8 zero-parameter route as canonical and demoting $m_c$ to a derived quantity (Section~\ref{sec:ledger}).
\item It registers \textbf{OPEN-FP-3-CKM} (the quark mixing matrix and quark $CP$ phase have no derivation in the CPP corpus), to be entered in the frontier registry at ship time via a flagged integration patch (Section~\ref{sec:ckm}).
\item It identifies the SF-line consistency threads SF-3$\leftrightarrow$SF-1, SF-3$\leftrightarrow$SF-2, and SF-3$\leftrightarrow$SF-4 that SF-3 unlocks for the SF-7 $\S10$ consistency-theorem family (Section~\ref{sec:discussion}).
\end{itemize}

% ============================================================
% SECTION 2: Substrate and the Four Bonded Shells
% ============================================================
\section{Substrate and the Four Bonded Shells}\label{sec:substrate}

Fix a reference vertex of the 600-cell and sort the remaining vertices by graph distance. The bonded distance shells produce a small set of regular sub-polytopes; the quark cages occupy four of them, with vertex counts
\[
V \in \{4,\,12,\,20,\,30\},
\]
identified respectively with the tetrahedral, icosahedral, dodecahedral, and icosidodecahedral shells. These four bonded shells are the geometric carriers of the four heavy quark flavours (strange, charm, bottom, top); the light up/down quarks belong to the lowest cage structures treated in the SM mass-generation papers and are not the subject of the heavy-spectrum formula below.

Colour structure enters through the Strong-Sector series. $SU(3)$ colour arises from cage-face permutation symmetry (SS-1), and the fundamental-representation Casimir
\[
\CF = \frac{N^2-1}{2N}\Big|_{N=3} = \frac{4}{3}
\]
is derived independently in SS-2. The combination $z\CF = 12\cdot\tfrac43 = 16$ supplies the top-quark relay multiplier in Section~\ref{sec:masses}. The single dimensionful calibration of the entire construction is the electron mass $\me$; everything else is geometry ($z$, $\phig$, the shell counts $V$) and the colour Casimir.

% ============================================================
% SECTION 3: The Zero-Parameter Quark Mass Spectrum
% ============================================================
\section{The Zero-Parameter Quark Mass Spectrum (Route A)}\label{sec:masses}

The heavy-quark masses follow the cage formula established in SM-8, with the exponent $7/3$ and the prefactor $\Mzero$ both \emph{derived} in SM-9 (pair count $\times$ linear cage dimension):
\begin{equation}
M_q = \me\,\frac{z}{\phig}\,V^{7/3}\quad(q=s,c,b),
\qquad
M_t = \me\,\frac{z}{\phig}\,V_t^{7/3}\times z\,\CF,
\label{eq:massformula}
\end{equation}
with the mass anchor
\begin{equation}
\Mzero \equiv \me\,\frac{z}{\phig} = 0.5110\;\mathrm{MeV}\times\frac{12}{1.6180} = 3.79\;\mathrm{MeV},
\label{eq:M0}
\end{equation}
$z=12$, $\CF=4/3$ (SS-2), and the top relay multiplier $z\CF = 16$. The four bonded shells of Section~\ref{sec:substrate} give the cage assignments and the predicted masses in Table~\ref{tab:masses}.

\begin{table}[H]
\centering
\caption{Zero-parameter heavy-quark mass spectrum from Eq.~\eqref{eq:massformula}. Inputs: $\me$, $z=12$, $\phig$, $\CF=4/3$. No shape parameters fitted. PDG values are pole/$\overline{\mathrm{MS}}$ reference masses; agreement is RMS $2.1\%$ across four orders of magnitude.}
\label{tab:masses}
\begin{tabular}{@{}llrrrrr@{}}
\toprule
Quark & Shell & $V$ & multiplier & CPP mass & PDG & error \\
\midrule
strange & tetrahedron       & 4  & 1            & $96.3$~MeV    & $93.4$~MeV    & $+3.1\%$ \\
charm   & icosahedron       & 12 & 1            & $1249$~MeV    & $1270$~MeV    & $-1.6\%$ \\
bottom  & dodecahedron      & 20 & 1            & $4115$~MeV    & $4180$~MeV    & $-1.6\%$ \\
top     & icosidodecahedron & 30 & $z\CF=16$    & $169{,}570$~MeV & $172{,}760$~MeV & $-1.8\%$ \\
\bottomrule
\end{tabular}
\end{table}

The root-mean-square residual is $2.1\%$, obtained from $\me$, the 600-cell geometry, and $SU(3)$ colour alone, with no parameters fitted. The single most consequential feature for the calibration question (Section~\ref{sec:ledger}) is that Eq.~\eqref{eq:massformula} \emph{predicts} the charm mass, $m_c = 1249$~MeV at $-1.6\%$, rather than taking it as an input. Charm is therefore a derived quantity in this route, and the spectrum carries a single dimensionful calibration, $\me$. The quoted residuals are computed against the stated PDG reference scheme (pole/$\overline{\mathrm{MS}}$); the scheme choice affects the empirical comparison, not the calibration --- no PDG mass is used to fix any parameter.

\begin{remark}
SM-10 supplies a finite-element chain-network \emph{mechanism} for the $7/3$ scaling; it is currently a calibrated geometric model (four fitted parameters against four data) pending a first-principles GPU closure. The mass \emph{values} in Table~\ref{tab:masses} do not depend on SM-10: they come from SM-8/SM-9. The SM-10 status is therefore a mechanism-depth caveat, not a mass-prediction caveat, and is carried in the conditional accounting (Section~\ref{sec:falsifiers}).
\end{remark}

% ============================================================
% SECTION 4: Strong Coupling and Electroweak-Strong Complementarity
% ============================================================
\section{Strong Coupling and Electroweak--Strong Mode Complementarity}\label{sec:alphas}

The strong coupling is the \emph{face-mode} fraction of the same 600-cell adjacency spectral trace that yields the Weinberg angle as its \emph{edge-mode} fraction. The adjacency operator $A$ is the $120\times120$ matrix encoding the vertex--vertex connectivity of the 600-cell. With $A$ so defined,
\begin{equation}
\alphas = \frac{1}{\phig}\cdot\frac{\tfrac13\Tr(A^3)}{\Tr(A^2)+\tfrac13\Tr(A^3)}
= \frac{1}{\phig}\cdot\frac{2400}{3840}
= \frac{5}{8\phig}\approx 0.386,
\label{eq:alphas}
\end{equation}
which matches the PDG running coupling at the charm scale, $\alphas(m_c)\approx 0.35$--$0.40$. The Weinberg angle inherited from SM-6 is the complementary edge-mode fraction,
\begin{equation}
\sinsqthetaW = \frac{3}{8\phig}\approx 0.232,
\end{equation}
and the two mode fractions are exactly complementary:
\begin{equation}
\boxed{\;\sinsqthetaW + \alphas = \frac{3}{8\phig} + \frac{5}{8\phig} = \frac{1}{\phig}\approx 0.618\;}
\label{eq:complementarity}
\end{equation}
The bare (pre-$\eta$) partition is $3/8 + 5/8 = 1$, corresponding to the unweighted adjacency contributions before propagation-efficiency ($\eta = 1/\phig$) scaling, and the topological ratio is
\begin{equation}
\frac{\alphas}{\sinsqthetaW} = \frac{F}{E} = \frac{1200}{720} = \frac{5}{3},
\end{equation}
the ratio of the 600-cell face and edge counts. \emph{Within the inherited SM-6/SM-7 framework, a single spectral trace supplies both couplings as complementary edge- and face-mode fractions} --- a mode-fraction correspondence, not a dynamical unification. This substrate-level electroweak--strong \emph{mode complementarity} is the strongest single $\S10$ consistency thread linking SF-3 to SF-1, SF-2, and SF-5: the same adjacency spectrum, partitioned into edge and face modes weighted by $\eta = 1/\phig$, sources both couplings with no independent calibration. Within SM-6/SM-7 these are \emph{structural correspondences} --- mode fractions of the adjacency spectrum that numerically match the measured couplings --- not gauge couplings obtained from renormalization-group running.

% ============================================================
% SECTION 5: The Quark Koide Phase
% ============================================================
\section{The Quark Koide Phase}\label{sec:koide}

The quark Koide phase derives from the same isotropic-self-energy-shift machinery that fixes the lepton Koide phase (SF-1), with a sector-appropriate strong-colour contribution added to the electroweak contribution. The base value is the $\Kthree$ eigenvalue ratio,
\begin{equation}
\cos\theta_0 = -K = -\frac{2}{3}.
\end{equation}
Quarks carry colour, so their $\Kthree$ cage faces feel the strong coupling on all $z=12$ nearest-neighbour bonds, producing a negative isotropic shift, alongside the same electroweak shift the leptons feel:
\begin{equation}
\varepsilon_S = -\frac{z\,\alphas}{z+1} = -\frac{60}{104\,\phig},
\qquad
\varepsilon_{EW} = +\frac{3}{52\,\phig},
\qquad
\varepsilon = \varepsilon_S + \varepsilon_{EW} = -\frac{27}{52\,\phig}.
\label{eq:shift}
\end{equation}
The shifted phase is
\begin{equation}
\cos\theta_{\mathrm{quark}}
= -\frac{2}{3}\left(1+\frac{\varepsilon}{2}\right)
= -\frac{2}{3}\left(1-\frac{27}{104\,\phig}\right)
\approx -0.5597
\;\Longrightarrow\;
\boxed{\;\theta_{\mathrm{quark}} = 124.04\dgr\;}
\label{eq:koide}
\end{equation}
versus the PDG-quark-mass-extracted value $124.09\dgr$ --- a $0.05\%$ agreement and a zero-parameter prediction.

\begin{proposition}[Phase--mass separation --- a bookkeeping observation]\label{prop:phaseindep}
This is a bookkeeping separation of inherited SM-7 structure, not a new derivation. In the retained SM-7 isotropic-shift formula, the quark Koide phase $\theta_{\mathrm{quark}}$ of Eq.~\eqref{eq:koide} is a function of $\alphas$, $\sinsqthetaW$, and $z$ only; it does not depend on the mass amplitude or on the charm mass $m_c$.
\end{proposition}

\noindent\textit{Consequence and SM-7 provenance.} An earlier reading (the SF-3 outline) suggested that the quark-phase machinery would need to be ``re-grounded'' on the derived $m_c$ once $m_c$ was demoted from calibration to prediction. Proposition~\ref{prop:phaseindep} shows this step is unnecessary. In SM-7, $\alphas = 5/(8\phig)$ and the quark-phase shift were obtained from the 600-cell adjacency spectrum alone --- not from PDG-extracted quark masses --- so no indirect $m_c$ dependence enters through a fit to the empirical phase. The charm mass entered the SM-7 construction only as the mass-\emph{amplitude} calibration $A_q$; because $A_q$ is an overall scale in the Koide parametrisation, it cancels from the ratio that fixes $\cos\theta_{\mathrm{quark}}$ and never appears in Eq.~\eqref{eq:shift} (the cancellation is written out explicitly in Appendix~\ref{app:shift}). This independence is conditional on one point worth stating explicitly: it holds because $\alphas$ is taken as the \emph{structural} value $5/(8\phig)$; were $\alphas$ instead extracted from a running-coupling fit at the charm scale, $m_c$ would re-enter indirectly through that fit. SF-3 adopts the structural value, so the independence stands. To be precise about scope: it is the \emph{retained SM-7 isotropic-shift formula} that carries no $m_c$ dependence; the phase therefore stands whichever mass route is used, and the SF-3 synthesis is correspondingly simpler than a hybrid: masses come entirely from the SM-8 cage formula, while $\alphas$ and $\theta_{\mathrm{quark}}$ come as derived structural observables independent of any mass calibration.

% ============================================================
% SECTION 6: Three Generations
% ============================================================
\section{Three Generations and the No-Fourth-Quark Prediction}\label{sec:generations}

The four bonded shells $V\in\{4,12,20,30\}$ exhibit a palindrome symmetry in the full 600-cell shell sequence, and antipodal identification in the tessellated lattice limits the Standard Model to exactly three generations (SM-8). The generation count is therefore \emph{selected within the SM-8 antipodal-identification model}, not assumed, and the same identification predicts \textbf{no fourth quark}. This is a zero-parameter structural prediction and a clean falsifier: discovery of a fourth-generation quark would directly contradict the antipodal-identification mechanism.

% ============================================================
% SECTION 7: Calibration Ledger
% ============================================================
\section{Calibration Ledger: the Route Adjudication}\label{sec:ledger}

Two shipped routes produce the heavy-quark masses. Adjudicating between them is the central drafting decision of SF-3, and it is settled here.

\paragraph{Route A (SM-8) --- zero-parameter cage formula.} Eq.~\eqref{eq:massformula} predicts all four heavy masses to RMS $2.1\%$ using only $\me$ + geometry + $\CF$. Crucially, $m_c$ is a \emph{prediction} ($1249$~MeV, $-1.6\%$), not an input. \textbf{One calibration ($\me$).}

\paragraph{Route B (SM-7) --- Koide $+$ $m_c$ calibration.} Uses the charm mass as a calibration constant to place the quark Koide phase, giving slightly better $b/t$ residuals ($m_b=4.24$~GeV, $1.4\%$; $m_t=169.8$~GeV, $1.7\%$) \emph{plus} the quark Koide phase and $\alphas = 5/(8\phig)$. \textbf{Two calibrations ($\me$, $m_c$).}

\paragraph{The adjudication.} Because Route A \emph{derives} $m_c$ from $\me$ at zero parameters, the $m_c$ calibration in Route B is \emph{redundant}: $m_c$ can be supplied by the Route-A prediction rather than fitted. The canonical SF-3 synthesis is therefore:
\begin{quote}
Adopt \textbf{Route A} as the canonical mass-spectrum route (all four heavy masses from $\me$ alone, RMS $2.1\%$, single calibration); \textbf{demote $m_c$ from calibration to derived quantity}. Retain Route B's \emph{structural} content --- $\alphas = 5/(8\phig)$, the complementarity $\sinsqthetaW+\alphas = 1/\phig$, and the quark Koide phase $124.04\dgr$ --- as the electroweak--strong complementarity layer. By Proposition~\ref{prop:phaseindep}, no re-grounding of the phase on a derived $m_c$ is needed: in the retained isotropic-shift formula, the phase does not depend on $m_c$.
\end{quote}

This restores the single-$\me$ calibration headline for the quark sector at the cost of slightly larger mass residuals than Route B's better $b/t$ ($\sim 0.3\%$). The trade is clearly worth it: a uniform single-calibration claim across charged leptons (SF-1), neutrinos (SF-4), and quarks (SF-3) is the spine of the SF-7 ``hierarchy without hierarchy'' argument, and a sub-percent accuracy gain is not worth a second calibration that contradicts that spine. With Route A adopted, the SF-7 master comparison table reads \textbf{one calibration ($\me$)} for all fermion masses. The remaining caveats --- SF-2's open $\eta$-dilution for boson \emph{absolute} masses, and the $\CKM$ gap below --- are not fermion-mass-calibration issues and do not disturb the single-$\me$ claim. In short: \textbf{Route A is now canonical; Route B's earlier $m_c$ calibration is superseded.} For completeness, the exponent $7/3$, the anchor $\Mzero$, and the top relay factor $z\CF$ are inherited structural outputs of SM-9 and SS-2; they are stated here as inherited, not re-tuned or re-argued in SF-3.

\begin{mdframed}
\textbf{Calibration adjudication, in one line.} \emph{Masses} (s,c,b,t) from the SM-8 cage formula on a \emph{single} calibration $\me$ ($m_c$ derived); \emph{$\alphas$ and the quark Koide phase} as inherited structural observables independent of any mass calibration. \textbf{Route A canonical; Route B's $m_c$ calibration superseded.} Zero shape parameters; $\CKM$ inherited-open.
\end{mdframed}

% ============================================================
% SECTION 8: CKM - Inherited Open
% ============================================================
\section{CKM Mixing --- Inherited Open}\label{sec:ckm}

There is \textbf{no $\CKM$-matrix derivation anywhere in the quark corpus}. SM-10, despite an incidental mention, is the finite-element scaling-mechanism paper, not a mixing paper. SF-3 therefore predicts the quark \emph{masses} and the generation count at zero parameters but does \emph{not} derive the $\CKM$ mixing angles or the quark $CP$ phase. We register this honestly:
\begin{openproblem}[OPEN-FP-3-CKM]
The quark mixing matrix $\CKM$ and the quark $CP$-violating phase are not derived in the CPP corpus. SF-3 inherits this as an open problem, to be entered in the frontier registry at ship time.
\end{openproblem}

\noindent This is structurally parallel to SF-4's open neutrino $\delta_{CP}$ in the limited sense that both flagships derive masses while deferring a mixing-sector $CP$ observable --- a parallel of posture, not an equivalence of difficulty or status. SF-4 ships seven of eight neutrino parameters with $\delta_{CP}$ deferred; SF-3 ships the quark masses and generation count with $\CKM$ deferred. Stating the parallel explicitly strengthens the SF-line's uniform ``masses derived, mixing-sector open'' posture. A candidate closure route --- to flag, not to pursue here --- is a quark-sector cage-mixing structure analogous to SM-5's $\Kthree \to$ tri-bimaximal derivation of the $\PMNS$ matrix; the $\CKM$ analog of that spectral-alignment mechanism is presently unscoped.

\begin{remark}
The quark $CP$ phase inside $\CKM$ is a \emph{separate} object from the neutrino $\delta_{CP}$ being pursued in the dedicated $\delta_{CP}$ window. SF-3's mixing gap therefore carries low cross-window collision risk: it is not that window's territory.
\end{remark}

% ============================================================
% SECTION 9: Conditional Accounting and Falsifiers
% ============================================================
\section{Conditional Accounting and Falsifiers}\label{sec:falsifiers}

\paragraph{Calibration ledger.} One dimensionful constant, $\me$. \textbf{Zero shape parameters.} The masses $m_s, m_c, m_b, m_t$ are all derived (SM-8/SM-9); $\alphas$ and $\theta_{\mathrm{quark}}$ are derived structural observables (SM-7); the generation count is selected within the SM-8 antipodal-identification model. The MeV scale enters only through the identification of the electron cage with the physical electron mass --- the same single calibration shared across SF-1, SF-3, and SF-4 --- not as an additional scale-setting step.

\paragraph{Inherited-open / conditional.}
\begin{itemize}[leftmargin=2em]
\item \textbf{OPEN-FP-3-CKM:} quark mixing and quark $CP$ phase undelivered (Section~\ref{sec:ckm}).
\item SM-10's first-principles cascade derivation remains a calibrated model (four parameters / four data) pending GPU closure. The mass \emph{values} do not depend on it.
\item Accuracy honesty: RMS $2.1\%$ on the masses under the single-$\me$ route; Route B gave marginally better $b/t$ at the cost of a second calibration (Section~\ref{sec:ledger}).
\end{itemize}

\paragraph{Falsifiers.}
\begin{enumerate}[leftmargin=2em]
\item \textbf{Fourth-generation quark.} Discovery of a fourth-generation quark contradicts the antipodal-identification three-generation prediction (Section~\ref{sec:generations}).
\item \textbf{Quark Koide phase.} A measured quark Koide phase departing materially from $124.04\dgr$ (well outside the present $0.05\%$ agreement) falsifies the isotropic-shift derivation (Section~\ref{sec:koide}).
\item \textbf{Mode complementarity.} A demonstration that $\sinsqthetaW$ and $\alphas$ do not share the single 600-cell spectral-trace origin --- e.g.\ an independent measurement forcing $\sinsqthetaW + \alphas \neq 1/\phig$ at the substrate scale --- falsifies Eq.~\eqref{eq:complementarity}.
\item \textbf{Mass-formula structure.} A heavy-quark mass found to violate the $V^{7/3}$ cage scaling beyond the stated residual band falsifies Eq.~\eqref{eq:massformula}.
\end{enumerate}

% ============================================================
% SECTION 10: Physical Interpretation
% ============================================================
\section{Physical Interpretation}\label{sec:interp}

In the CPP ontology, a quark is a colour-charged cage of Conscious Points occupying one of the four bonded distance-shells of the 600-cell. The cage's vertex count $V$ sets its mass through the $V^{7/3}$ scaling of Eq.~\eqref{eq:massformula}: the exponent $7/3$ counts displacement-increment pair interactions times the linear cage dimension (SM-9), so a larger cage stores proportionally more $\SSV$ field energy. The top quark's extra factor $z\CF=16$ reflects the colour-Casimir-weighted coordination of the icosidodecahedral shell relaying its self-energy across all $z=12$ neighbours. The strong coupling and the Weinberg angle are not separate inputs but two harmonics of one substrate vibration: $\alphas$ is the share of the adjacency spectral trace carried by face ($A^3$) modes and $\sinsqthetaW$ the share carried by edge ($A^2$) modes, their sum fixed at $1/\phig$ by the propagation-efficiency factor. ($A^3$ counts oriented face-adjacency walks, which correspond to colour-flux relays across cage faces; $A^2$ counts edge round-trips.) The quark Koide phase is the geometric angle by which colour and electroweak self-energy shifts rotate the $\Kthree$ eigenphase away from its bare $-2/3$ value.

\subsection{CP/GP Signature at This Scale}\label{sec:signature}

\textbf{Load-bearing axioms.} The quark-sector derivation rests on the 600-cell topology axiom (supplying $z=12$ and the shell counts $V$), the propagation-efficiency axiom (supplying the $1/\phig$ weighting common to $\alphas$, $\sinsqthetaW$, and $\Mzero$), the cage-volume scaling axiom (supplying $\Mzero = \me z/\phig$ and the $7/3$ exponent through SM-9), and the lattice-scale grounding axiom (fixing MeV units). $SU(3)$ colour enters through the Strong-Sector cage-face-permutation axiom and the derived Casimir $\CF=4/3$ (SS-2).

\textbf{Visible-versus-smoothed discreteness.} The lattice discreteness is \emph{visible} in this paper as the integer shell counts $V\in\{4,12,20,30\}$, the integer edge/face counts $E=720$, $F=1200$ in the topological ratio $\alphas/\sinsqthetaW = 5/3$, the coordination $z=12$, and the $\phig$ factors throughout. It is \emph{smoothed} only in the running-coupling comparison, where $\alphas = 5/(8\phig)$ is matched to a renormalization-scheme-dependent effective value averaged over many substrate interactions at the charm scale.

\textbf{Macroscopic shadow.} The empirical regularity that the quark mass spectrum spans orders of magnitude yet sits on a single geometric progression --- and that the strong and electroweak couplings sum to a fixed golden-ratio constant --- is, within the CPP ontology, the macroscopic shadow of the sub-quantum claim that all four heavy quarks are cages on one lattice and both couplings are mode fractions of one adjacency spectrum. What conventional physics records as four unrelated quark masses plus two unrelated couplings is, in CPP, one cage-scaling law plus one spectral partition.

% ============================================================
% SECTION 11: CPP-to-Conventional-Physics Mapping
% ============================================================
\section{CPP-to-Conventional-Physics Mapping}\label{sec:mapping}

\begin{table}[H]
\centering
\caption{Structural correspondence between the CPP quark-sector mechanism and the conventional Standard-Model description. The mapping is at the level of degrees of freedom, symmetry, and observable, not a literal identity --- CPP operates at finer granularity than perturbative QFT.}
\label{tab:mapping}
\begin{tabular}{@{}p{0.46\textwidth}p{0.46\textwidth}@{}}
\toprule
\textbf{CPP mechanism} & \textbf{Conventional physics} \\
\midrule
Colour-charged CP cage on a bonded shell, vertex count $V$ & A quark flavour with a definite mass \\
$V^{7/3}$ cage-energy scaling with anchor $\Mzero=\me z/\phig$ & The measured heavy-quark mass hierarchy \\
Face-mode fraction of the adjacency spectral trace & The strong coupling $\alphas$ \\
Edge-mode fraction of the same spectral trace & The Weinberg angle $\sinsqthetaW$ \\
Isotropic colour$+$electroweak self-energy shift on $\Kthree$ & The quark Koide phase $\theta_{\mathrm{quark}}$ \\
Antipodal identification in the tessellated lattice & Exactly three generations; no fourth quark \\
(undelivered) cage-mixing alignment structure & The $\CKM$ matrix and quark $CP$ phase \\
\bottomrule
\end{tabular}
\end{table}

% ============================================================
% SECTION 12: Discussion - SF-line placement
% ============================================================
\section{Discussion: SF-line Placement and the $\S10$ Consistency Threads}\label{sec:discussion}

SF-3 is more than a quark mass table: it supplies three of the cross-sector consistency threads that the SF-7 grand-unification paper assembles into its $\S10$ consistency-theorem family. Each follows the three-clause consistency template (single shared source; no independent calibration; no conflicting prediction).

\begin{itemize}[leftmargin=2em]
\item \textbf{SF-3 $\leftrightarrow$ SF-2.} The mode-fraction structure: $\sinsqthetaW = 3/(8\phig)$ (SF-2) and $\alphas = 5/(8\phig)$ (SF-3) come from \emph{one} spectral trace, complementary to $1/\phig$. A strong consistency clause: a single spectral source, no independent calibration.
\item \textbf{SF-3 $\leftrightarrow$ SF-1.} The Koide-phase derivation machinery (isotropic self-energy shift on the $\Kthree$ face) and the $\Mzero=\me z/\phig$ anchor are shared. Consistency clause: lepton and quark phases use the same mechanism with sector-appropriate shifts (electroweak only for leptons; electroweak $+$ strong for quarks).
\item \textbf{SF-3 $\leftrightarrow$ SF-4.} The $\Mzero$ anchor and four-cage taxonomy are shared, as is the \emph{masses-derived / mixing-open} parallel ($\CKM$ open in SF-3 $\leftrightarrow$ $\delta_{CP}$ open in SF-4). Consistency clause: both sectors derive masses from $\Mzero$ and defer their mixing-sector $CP$ structure, with no conflicting calibration.
\end{itemize}

These threads are why the single-$\me$ adjudication of Section~\ref{sec:ledger} matters beyond the quark sector: it lets the SF-7 master table commit to one calibration across all fermion masses, which is the load-bearing claim of the unification argument.

% ============================================================
% SECTION 13: Conclusion
% ============================================================
\section{Conclusion}\label{sec:conclusion}

The complete heavy-quark mass spectrum, the strong coupling, the quark Koide phase, and the three-generation count all follow from a single calibration ($\me$) plus 600-cell geometry and $SU(3)$ colour, with zero shape parameters. The earlier ambiguity over whether the quark sector needed a second calibration ($m_c$) is resolved in favour of a single calibration by adopting the SM-8 zero-parameter route as canonical and demoting $m_c$ to a derived quantity; a worked phase--mass separation (Proposition~\ref{prop:phaseindep}) shows that, in the retained SM-7 isotropic-shift formula, the quark Koide phase does not depend on $m_c$. The quark mixing matrix remains honestly undelivered (OPEN-FP-3-CKM), the structural analog of SF-4's open $\delta_{CP}$. The mass values do not depend on SM-10's still-calibrated chain-network mechanism; its first-principles closure is an open mechanism-depth question, not a caveat on the mass predictions.

\subsection{Swarm-Validation Contribution}\label{sec:swarm}

SF-3 contributes the following zero-parameter empirical correspondences to the cumulative CPP swarm, all from the same unchanged axiom stack: the four heavy-quark masses (RMS $2.1\%$), the strong coupling $\alphas = 5/(8\phig)$, the exact electroweak--strong complementarity $\sinsqthetaW+\alphas = 1/\phig$, the quark Koide phase $124.04\dgr$ ($0.05\%$), and the three-generation count (selected within the SM-8 antipodal-identification model). Because each correspondence is fixed by geometry once $\me$ is set, the probability of reproducing them jointly by accident or hidden tuning scales as $(\text{residual band}/\text{typical parameter space})^{N}$ in the number of independent correspondences $N$, and is already very small at the swarm's present size. (The cumulative programme counter is updated in \texttt{predictions.md} at ship time via a flagged integration patch.)

\subsection{Problem Status After This Paper}\label{sec:status}

\begin{itemize}[leftmargin=2em]
\item Single-$\me$ quark-calibration question: \textbf{RESOLVED} (Route A canonical; $m_c$ demoted to derived).
\item \textbf{NEW:} OPEN-FP-3-CKM registered (quark mixing matrix and quark $CP$ phase undelivered).
\item SF-7 $\S10$ threads SF-3$\leftrightarrow$SF-1, SF-3$\leftrightarrow$SF-2, SF-3$\leftrightarrow$SF-4 identified and ready for the consistency-theorem family.
\end{itemize}

% ============================================================
% ACKNOWLEDGEMENTS
% ============================================================
\section*{Acknowledgements}

This paper is a synthesis and reframing of shipped CPP results (SM-7, SM-8, SM-9, SM-10, SS-1, SS-2, SM-6) into a unified quark-sector account; it introduces no new derivation. Thomas Lee Abshier (ND) conceived the Conscious Point Physics program, the 600-cell substrate hypothesis, and the physical intuitions underlying the cage-mass and isotropic-shift mechanisms. Claude Opus (Anthropic) performed the corpus survey, the calibration-route adjudication, and the drafting, and identified and documented the phase--mass bookkeeping separation (Proposition~\ref{prop:phaseindep}, with the explicit $A_q$ cancellation given in Appendix~\ref{app:shift}). The multi-AI review panel (ChatGPT, Grok, Copilot) reviewed the paper across four rounds; the panel's sharpest contribution was the observation, incorporated in Section~\ref{sec:koide}, that the phase--mass independence holds only because $\alphas$ is taken as the structural value rather than a charm-scale running-coupling fit. The 600-cell lattice hypothesis (AXIM-2) remains an axiom. Repository: \url{https://github.com/Hyperphysics-Institute/CPP}; OSF: \url{https://osf.io/9dfya/} (DOI 10.17605/OSF.IO/JXE8D).

\appendix

% ============================================================
% APPENDIX A: SM-7 isotropic-shift formula and the A_q cancellation
% ============================================================
\section{The SM-7 Isotropic-Shift Formula and the $A_q$ Cancellation}\label{app:shift}

This appendix makes Proposition~\ref{prop:phaseindep} self-contained by writing out the retained SM-7 isotropic-shift formula and showing explicitly that the mass amplitude $A_q$ --- the only channel through which the charm mass $m_c$ entered the SM-7 construction --- cancels from the equation that fixes the quark Koide phase. Nothing here is new to SF-3; it is a transcription of inherited SM-7 structure, included so the bookkeeping separation can be checked directly rather than taken on trust.

\paragraph{The Koide parametrisation.} Within a generation triplet, the construction writes the three square-root masses as
\begin{equation}
\sqrt{m_i} = A_q\left(1 + \sqrt{2}\,\cos\!\left(\theta + \frac{2\pi i}{3}\right)\right), \qquad i = 1,2,3,
\label{eq:koidepar}
\end{equation}
with a single overall amplitude $A_q$ and a single phase $\theta$. The amplitude $A_q$ sets the absolute mass scale of the triplet; the phase $\theta$ sets the \emph{relative} spacing of the three masses. (The Koide ratio $K = (\sum_i m_i)/(\sum_i \sqrt{m_i})^2 = 2/3$ is recovered identically for any $A_q,\theta$, since numerator and denominator both scale as $A_q^2$.)

\paragraph{The isotropic shift acts on $\theta$, not $A_q$.} The bare $\Kthree$ eigenphase is $\cos\theta_0 = -K = -2/3$. The colour-plus-electroweak isotropic self-energy shift $\varepsilon$ of Eq.~\eqref{eq:shift} perturbs the eigenphase multiplicatively,
\begin{equation}
\cos\theta_{\mathrm{quark}} = -\frac{2}{3}\left(1 + \frac{\varepsilon}{2}\right), \qquad \varepsilon = \varepsilon_S + \varepsilon_{EW} = -\frac{z\,\alphas}{z+1} + \frac{3}{52\,\phig} = -\frac{27}{52\,\phig},
\label{eq:shiftapp}
\end{equation}
and $\varepsilon$ is built from $\{\alphas,\ \sinsqthetaW,\ z\}$ alone. The amplitude $A_q$ does not appear in Eq.~\eqref{eq:shiftapp}.

\paragraph{The cancellation.} Rescale the amplitude $A_q \to \lambda A_q$ for any $\lambda > 0$. By Eq.~\eqref{eq:koidepar} every $\sqrt{m_i}$ scales by $\lambda$, so every mass scales by $\lambda^2$ and every \emph{ratio} $m_i/m_j$ is unchanged. The phase $\theta_{\mathrm{quark}}$ is fixed by the mass ratios --- equivalently by Eq.~\eqref{eq:shiftapp}, which contains no $A_q$ --- so it is invariant under the rescaling: $\partial\theta_{\mathrm{quark}}/\partial A_q = 0$. Because $m_c$ entered the SM-7 construction \emph{only} as the calibration that fixed $A_q$ for the quark triplet, it follows that $\partial\theta_{\mathrm{quark}}/\partial m_c = 0$. This is the content of Proposition~\ref{prop:phaseindep}.

\paragraph{The one conditional.} The argument uses $\alphas = 5/(8\phig)$ as a fixed structural input. If instead $\alphas$ were taken from a running-coupling determination at the charm scale $\mu \sim m_c$, then $\alphas = \alphas(m_c)$ would carry an implicit $m_c$ dependence into $\varepsilon$, and hence into Eq.~\eqref{eq:shiftapp}, and the independence would fail. SF-3 adopts the structural value throughout, so $\theta_{\mathrm{quark}}$ is genuinely $m_c$-independent; the demotion of $m_c$ from calibration to derived quantity (Section~\ref{sec:ledger}) therefore leaves the quark Koide phase untouched, with no re-grounding required.

% ============================================================
% BIBLIOGRAPHY (inline; master entry abshier2026sf3 added at ship)
% ============================================================
\begin{thebibliography}{99}

\bibitem{abshier_sm6} T.~L.~Abshier et al., \emph{The Weinberg Angle from 600-Cell Spectral Traces}, SM-6, Hyperphysics Institute (2026).
\bibitem{abshier_sm7} T.~L.~Abshier et al., \emph{Strong Coupling, Mode Complementarity, and the Quark Koide Phase}, SM-7, Hyperphysics Institute (2026).
\bibitem{abshier_sm8} T.~L.~Abshier et al., \emph{Quark Generations and the Zero-Parameter Cage Mass Formula}, SM-8 v4.0, Hyperphysics Institute (2026).
\bibitem{abshier_sm9} T.~L.~Abshier et al., \emph{Derivation of the $7/3$ Cage-Mass Exponent and the Anchor $M_0$}, SM-9, Hyperphysics Institute (2026).
\bibitem{abshier_sm10} T.~L.~Abshier et al., \emph{Finite-Element Chain-Network Mechanism for Cage-Mass Scaling}, SM-10, Hyperphysics Institute (2026).
\bibitem{abshier_ss1} T.~L.~Abshier et al., \emph{$SU(3)$ Colour from Cage-Face Permutation Symmetry}, SS-1, Hyperphysics Institute (2026).
\bibitem{abshier_ss2} T.~L.~Abshier et al., \emph{The Colour Casimir $C_F=4/3$ from the 600-Cell}, SS-2, Hyperphysics Institute (2026).
\bibitem{abshier_sf2} T.~L.~Abshier and {Claude Opus}, \emph{SF-2: Electroweak Cage-Boson Unification from 600-Cell Geometry}, v1.01, Hyperphysics Institute (2026).
\bibitem{abshier_sf4} T.~L.~Abshier et al., \emph{SF-4: The Neutrino Sector from 600-Cell Geometry}, v4.4, Hyperphysics Institute (2026).
\bibitem{koide1983} Y.~Koide, \emph{A fermion-boson composite model of quarks and leptons}, Phys.\ Lett.\ B \textbf{120}, 161 (1983).
\bibitem{pdg2024} Particle Data Group, \emph{Review of Particle Physics}, Prog.\ Theor.\ Exp.\ Phys.\ (2024).

\end{thebibliography}

\end{document}
